\begin{titlepage}
\centering
{\Large Notas de entrevista\par}
\vspace{1.2cm}
{\huge\bfseries \notetitle\par}
\vspace{0.8cm}
{\large \noteauthors\par}
\vspace{0.3cm}
{\large \notedate\par}
\vspace{1.2cm}
\ifdefempty{\videocoverpath}{}{\includegraphics[width=0.82\textwidth,height=0.45\textheight,keepaspectratio]{\videocoverpath}\par}
\vfill
\begin{tcolorbox}[width=0.9\textwidth, colback=black!2!white, colframe=black!60, sharp corners]
\textbf{Autor o canal del video}: \videochannel\par
\textbf{Fecha de publicación}: \videopublishdate\par
\textbf{Duración del video}: \videoduration\par
\ifdefempty{\videourl}{}{\textbf{Enlace del video}: \href{\videourl}{\nolinkurl{\videourl}}\par}
\textbf{Página del proyecto}: \href{\repourl}{\nolinkurl{\repourl}}\par
\end{tcolorbox}
\vspace{0.3cm}
{\small Esta charla se concentra en dos desplazamientos de enfoque en la infraestructura de los modelos grandes durante 2025: en la primera mitad del año, en torno a la clusterización de la inferencia y la separación PD; en la segunda mitad, en torno al scaling de la infraestructura de RL / Agent. Los invitados provienen tanto del ámbito académico como de la industria, por lo que la charla hace énfasis particular en «cómo el sistema absorbe nuevas workloads».\par}
\end{titlepage}
\tableofcontents
\newpage

